La separación ciega de fuentes es aquella que busca separar las señales originales a partir únicamente de las mezclas observadas, sin información previa sobre las fuentes ni sobre el proceso de mezcla.\cite{cardoso1998blind}

Dos técnicas clásicas son:
\begin{itemize}
    \item \textbf{\textit{Independent Component Analysis ICA (Análisis de componentes independientes)}:} busca componentes estadísticamente independientes, muy usada en aplicaciones de separación de voz.

    \item \textbf{\textit{Non-negative Matrix Factorization NMF (Factorización no negativa de matrices)}:} descompone el espectrograma en bases y activaciones, útil en música.
    
\end{itemize}
